\documentclass[xcolor=dvipsnames]{ctexbeamer}
\usetheme{Madrid}
\usecolortheme{default}
\usepackage{amsmath,amssymb,amsfonts}
\usepackage{graphicx}
\usepackage{xcolor}
\usepackage{hyperref}

\title{对称性与守恒律}
\subtitle{从拉格朗日到诺特定理}
\author{物理学演讲}
\date{\today}

\begin{document}

\frame{\titlepage}

\begin{frame}
\frametitle{目录}
\tableofcontents
\end{frame}

\section{引入问题}

\begin{frame}
\frametitle{分析力学的基本框架}
\begin{block}{拉格朗日方程}
$$\frac{d}{dt}\left(\frac{\partial L}{\partial \dot{q}_i}\right) - \frac{\partial L}{\partial q_i} = 0$$
\end{block}

\pause

\begin{itemize}
\item $L = T - V$ 是拉格朗日量（动能减势能）
\item $q_i$ 是广义坐标
\item $\dot{q}_i$ 是广义速度
\end{itemize}
\end{frame}

\begin{frame}
\frametitle{循环坐标}
\begin{definition}
如果拉格朗日量 $L$ 不显含某个坐标 $q_i$，即 $\frac{\partial L}{\partial q_i} = 0$，则称 $q_i$ 为\textbf{循环坐标}或\textbf{可遗坐标}。
\end{definition}

\pause

\begin{block}{重要结论}
当 $q_i$ 是循环坐标时，对应的广义动量守恒：
$$p_i = \frac{\partial L}{\partial \dot{q}_i} = \text{常数}$$
\end{block}
\end{frame}

\begin{frame}
\frametitle{核心问题}
\begin{center}
\Large \textcolor{red}{对称性的本质是什么？}
\end{center}

\pause

\begin{itemize}
\item 为什么"拉氏量不含某坐标"就导致守恒律？
\pause
\item 这背后是否有更深刻的原理？
\pause
\item 守恒律与对称性之间是否存在普遍联系？
\end{itemize}
\end{frame}

\section{守恒律的经典认识}

\begin{frame}
\frametitle{循环坐标与广义动量守恒}
\begin{example}[空间平移]
若 $L$ 不显含位置坐标 $x$：
$$\frac{\partial L}{\partial x} = 0 \quad \Rightarrow \quad p_x = \frac{\partial L}{\partial \dot{x}} = \text{常数}$$
\textbf{线动量守恒}
\end{example}

\pause

\begin{example}[旋转]
若 $L$ 不显含角坐标 $\theta$：
$$\frac{\partial L}{\partial \theta} = 0 \quad \Rightarrow \quad p_\theta = \frac{\partial L}{\partial \dot{\theta}} = \text{常数}$$
\textbf{角动量守恒}
\end{example}
\end{frame}

\begin{frame}
\frametitle{时间平移与能量守恒}
\begin{block}{能量函数}
定义能量函数（哈密顿量）：
$$H = \sum_i \dot{q}_i \frac{\partial L}{\partial \dot{q}_i} - L$$
\end{block}

\pause

\begin{theorem}
当拉格朗日量 $L$ 不显含时间 $t$，即 $\frac{\partial L}{\partial t} = 0$ 时，能量函数 $H$ 守恒：
$$\frac{dH}{dt} = 0$$
\end{theorem}
\end{frame}

\section{对称性的深层含义}

\begin{frame}
\frametitle{什么是"不变"？}
\begin{center}
\Large 对称性指什么"不变"？
\end{center}

\pause

\begin{enumerate}
\item \textbf{第一层}：拉氏量本身不变 \quad $L' = L$
\pause
\item \textbf{第二层}：拉氏量至多相差全导数 \quad $L' = L + \frac{df}{dt}$
\pause
\item \textbf{第三层}：\textcolor{red}{作用量的变分不变}（最根本） \quad $\delta S' = \delta S$
\end{enumerate}
\end{frame}

\begin{frame}
\frametitle{作用量的定义}
\begin{definition}[作用量]
$$S = \int_{t_1}^{t_2} L(q, \dot{q}, t) \, dt$$
\end{definition}

\pause

\begin{block}{哈密顿原理}
系统的真实运动使作用量取极值（通常是极小值）：
$$\delta S = 0$$
\end{block}

\pause

\textbf{关键洞察}：运动方程只依赖于作用量，而非拉氏量本身！
\end{frame}

\begin{frame}
\frametitle{全导数不影响运动方程}
\begin{theorem}
如果两个拉格朗日量相差一个全导数：
$$L' = L + \frac{df(q, t)}{dt}$$
则它们给出相同的运动方程。
\end{theorem}

\pause

\begin{proof}[证明思路]
$$S' = \int_{t_1}^{t_2} L' \, dt = \int_{t_1}^{t_2} L \, dt + \int_{t_1}^{t_2} \frac{df}{dt} \, dt$$
$$= S + [f(q(t_2), t_2) - f(q(t_1), t_1)]$$

由于端点固定，$\delta S' = \delta S$，故运动方程相同。
\end{proof}
\end{frame}

\begin{frame}
\frametitle{对称性的精确定义}
\begin{definition}[连续对称性]
如果存在连续变换：
$$t \to t' = t + \delta t, \quad q_i \to q'_i = q_i + \delta q_i$$
使得拉氏量至多相差全导数（即作用量至多相差边界项），则称系统具有该\textbf{连续对称性}。
\end{definition}

\pause

\begin{alertblock}{注意}
当端点固定时，边界项的变分为零，因此作用量的变分不变：$\delta S' = \delta S$。
\end{alertblock}
\end{frame}

\section{诺特定理}

\begin{frame}
\frametitle{Emmy Noether (1882-1935)}
\begin{columns}
\column{0.6\textwidth}
\begin{itemize}
\item 德国数学家
\item 1918年发表著名论文
\item 建立对称性与守恒律的普遍联系
\item 20世纪最伟大的数学家之一
\end{itemize}

\column{0.4\textwidth}
\begin{center}
\textit{[Emmy Noether的照片位置]}
\end{center}
\end{columns}
\end{frame}

\begin{frame}
\frametitle{诺特定理的表述}
\begin{theorem}[诺特定理]
\textbf{每一个连续对称性对应一个守恒律。}
\end{theorem}

\pause

更精确地说：

\begin{block}{诺特定理（数学形式）}
若系统在无穷小变换
$$\delta t = \epsilon \tau(q, t), \quad \delta q_i = \epsilon \xi_i(q, t)$$
下具有对称性，则存在守恒量：
$$Q = \sum_i p_i \xi_i - H \tau = \text{常数}$$
其中 $p_i = \frac{\partial L}{\partial \dot{q}_i}$，$H$ 是哈密顿量。
\end{block}
\end{frame}

\begin{frame}
\frametitle{对称性与守恒律的对应}
\begin{center}
\begin{tabular}{|c|c|}
\hline
\textbf{对称性} & \textbf{守恒律} \\
\hline
空间平移对称性 & 动量守恒 \\
\hline
空间旋转对称性 & 角动量守恒 \\
\hline
时间平移对称性 & 能量守恒 \\
\hline
规范对称性 & 电荷守恒 \\
\hline
\end{tabular}
\end{center}

\pause

\vspace{1em}
\begin{alertblock}{深刻意义}
守恒律不是巧合，而是\textcolor{red}{对称性的必然结果}！
\end{alertblock}
\end{frame}

\section{诺特定理前后的认识对比}

\begin{frame}
\frametitle{诺特定理之前 (17-19世纪)}
\begin{block}{已知的守恒律}
\begin{itemize}
\item 动量守恒（牛顿）
\item 能量守恒（焦耳、亥姆霍兹）
\item 角动量守恒（开普勒）
\end{itemize}
\end{block}

\pause

\begin{block}{拉格朗日的循环坐标}
\begin{itemize}
\item 18世纪末，拉格朗日注意到循环坐标与守恒量的关系
\item 但这只是\textbf{个案观察}，缺乏统一理解
\end{itemize}
\end{block}
\end{frame}

\begin{frame}
\frametitle{认识的局限性}
\textbf{19世纪的物理学家知道：}
\begin{itemize}
\item 空间均匀 $\Rightarrow$ 动量守恒
\item 空间各向同性 $\Rightarrow$ 角动量守恒
\item 时间均匀 $\Rightarrow$ 能量守恒
\end{itemize}

\pause

\vspace{1em}
\textbf{但他们不清楚：}
\begin{itemize}
\item 这些是\textbf{偶然}还是\textbf{必然}？
\item 背后是否有\textbf{统一的数学原理}？
\item 是否存在\textbf{其他未发现的守恒律}？
\end{itemize}
\end{frame}

\begin{frame}
\frametitle{诺特定理的革命性}
\begin{enumerate}
\item \textbf{统一性}
\begin{itemize}
\item 将所有守恒律纳入统一框架
\item 对称性是更基本的概念
\end{itemize}

\pause

\item \textbf{预言性}
\begin{itemize}
\item 可以通过寻找对称性来发现新守恒律
\item 例如：规范对称性 $\Rightarrow$ 电荷守恒
\end{itemize}

\pause

\item \textbf{深刻性}
\begin{itemize}
\item 揭示了物理定律与时空性质的深层联系
\item 成为现代物理学的基石
\end{itemize}
\end{enumerate}
\end{frame}

\section{具体例子}

\begin{frame}
\frametitle{例1：空间平移对称性}
\begin{example}[自由粒子]
$$L = \frac{1}{2}m(\dot{x}^2 + \dot{y}^2 + \dot{z}^2)$$
\end{example}

\pause

\textbf{对称性}：空间平移 $\vec{r} \to \vec{r} + \vec{a}$

\pause

拉氏量不含位置坐标 $x, y, z$：
$$\frac{\partial L}{\partial x} = \frac{\partial L}{\partial y} = \frac{\partial L}{\partial z} = 0$$

\pause

\textbf{守恒量}：
$$\vec{p} = \frac{\partial L}{\partial \dot{\vec{r}}} = m\dot{\vec{r}} = \text{常数}$$

\textcolor{blue}{动量守恒！}
\end{frame}

\begin{frame}
\frametitle{例2：旋转对称性}
\begin{example}[中心力场]
$$L = \frac{1}{2}m(\dot{r}^2 + r^2\dot{\theta}^2 + r^2\sin^2\theta \, \dot{\phi}^2) - V(r)$$
势能只依赖于 $r$
\end{example}

\pause

\textbf{对称性}：绕 $z$ 轴旋转，$\phi \to \phi + \delta\phi$

\pause

$$\frac{\partial L}{\partial \phi} = 0$$

\pause

\textbf{守恒量}：
$$p_\phi = \frac{\partial L}{\partial \dot{\phi}} = mr^2\sin^2\theta \, \dot{\phi} = L_z = \text{常数}$$

\textcolor{blue}{角动量分量守恒！}
\end{frame}

\begin{frame}
\frametitle{完整推导：能量守恒}
考虑时间平移 $t \to t + \epsilon$

\pause

拉氏量变化：
$$\delta L = \frac{\partial L}{\partial t} \epsilon$$

\pause

若 $\frac{\partial L}{\partial t} = 0$（时间平移对称性），则：

\pause

计算作用量的变化：
$$\frac{dH}{dt} = \sum_i \dot{q}_i \frac{dp_i}{dt} + \sum_i p_i \ddot{q}_i - \sum_i \frac{\partial L}{\partial q_i}\dot{q}_i - \sum_i \frac{\partial L}{\partial \dot{q}_i}\ddot{q}_i - \frac{\partial L}{\partial t}$$

\pause

利用拉格朗日方程和 $\frac{\partial L}{\partial t} = 0$：
$$\frac{dH}{dt} = 0$$

\textcolor{blue}{能量守恒！}
\end{frame}

\section{现代物理中的应用}

\begin{frame}
\frametitle{量子力学中的对称性}
\begin{itemize}
\item \textbf{对称性} $\Leftrightarrow$ \textbf{算符对易}
\pause
\item 例如：平移算符 $\hat{T}$ 与哈密顿算符 $\hat{H}$ 对易
$$[\hat{T}, \hat{H}] = 0 \quad \Rightarrow \quad \text{动量守恒}$$
\pause
\item 对称性决定量子数和选择定则
\end{itemize}
\end{frame}

\begin{frame}
\frametitle{规范对称性与电荷守恒}
\begin{block}{规范变换}
电磁场中，势函数的变换：
$$\vec{A} \to \vec{A} + \nabla\chi, \quad \phi \to \phi - \frac{\partial \chi}{\partial t}$$
不改变物理规律（规范对称性）
\end{block}

\pause

\begin{alertblock}{诺特定理的应用}
规范对称性 $\Rightarrow$ \textbf{电荷守恒}
$$\frac{\partial \rho}{\partial t} + \nabla \cdot \vec{j} = 0$$
\end{alertblock}
\end{frame}

\begin{frame}
\frametitle{粒子物理标准模型}
\begin{itemize}
\item 标准模型基于\textbf{规范对称性}：$SU(3) \times SU(2) \times U(1)$
\pause
\item 每个对称性对应一种相互作用：
\begin{itemize}
\item $SU(3)$：强相互作用
\item $SU(2) \times U(1)$：电弱相互作用
\end{itemize}
\pause
\item 对称性破缺导致质量产生（Higgs机制）
\end{itemize}

\vspace{1em}
\pause

\begin{center}
\Large \textcolor{red}{对称性是现代物理学的语言！}
\end{center}
\end{frame}

\section{总结}

\begin{frame}
\frametitle{总结}
\begin{enumerate}
\item \textbf{对称性是守恒律的根源}
\begin{itemize}
\item 不是守恒律导致对称性，而是对称性导致守恒律
\end{itemize}

\pause

\item \textbf{从经验到理论的飞跃}
\begin{itemize}
\item 诺特定理之前：零散的经验事实
\item 诺特定理之后：统一的理论框架
\end{itemize}

\pause

\item \textbf{20世纪物理学的基石}
\begin{itemize}
\item 量子场论
\item 粒子物理标准模型
\item 广义相对论
\end{itemize}
\end{enumerate}
\end{frame}

\begin{frame}
\begin{center}
{\Huge 谢谢！}

\vspace{2em}

{\Large 问题与讨论}
\end{center}
\end{frame}

\appendix

\section{附录A：点粒子力学中的诺特定理详细证明}

\begin{frame}
\frametitle{A.1 准备工作}
\begin{block}{拉格朗日量与作用量}
考虑一个具有 $n$ 个自由度的力学系统：
\begin{itemize}
\item 广义坐标：$q_i(t)$，$i = 1, 2, \ldots, n$
\item 拉格朗日量：$L(q_i, \dot{q}_i, t)$
\item 作用量：$S = \int_{t_1}^{t_2} L(q_i, \dot{q}_i, t) \, dt$
\end{itemize}
\end{block}

\pause

\begin{block}{欧拉-拉格朗日方程}
运动方程由变分原理 $\delta S = 0$ 得到：
$$\frac{d}{dt}\left(\frac{\partial L}{\partial \dot{q}_i}\right) - \frac{\partial L}{\partial q_i} = 0$$
\end{block}

\pause

\begin{block}{广义动量与哈密顿量}
$$p_i = \frac{\partial L}{\partial \dot{q}_i}, \quad H = \sum_i p_i \dot{q}_i - L$$
\end{block}
\end{frame}

\begin{frame}
\frametitle{A.2 无穷小对称变换}
\begin{block}{变换形式}
考虑依赖于无穷小参数 $\epsilon$ 的变换：
\begin{align*}
t &\to t' = t + \epsilon \tau(q, t) \\
q_i &\to q_i' = q_i + \epsilon \xi_i(q, t)
\end{align*}
其中 $\tau$ 和 $\xi_i$ 是给定的函数。
\end{block}

\pause

\begin{block}{坐标的全变分}
沿着轨迹，坐标的实际变化为：
$$\delta q_i = q_i'(t') - q_i(t) = \epsilon \xi_i + \dot{q}_i \epsilon \tau$$
这包含了两部分：
\begin{itemize}
\item $\epsilon \xi_i$：同一时刻的坐标改变
\item $\dot{q}_i \epsilon \tau$：时间改变引起的坐标改变
\end{itemize}
\end{block}
\end{frame}

\begin{frame}
\frametitle{A.3 速度的变换}
\begin{block}{速度变换公式}
速度 $\dot{q}_i = \frac{dq_i}{dt}$ 在变换下的变化：
$$\delta \dot{q}_i = \frac{d(\delta q_i)}{dt} - \dot{q}_i \frac{d(\epsilon \tau)}{dt}$$
\end{block}

\pause

\begin{block}{推导}
\begin{align*}
\delta \dot{q}_i &= \frac{d}{dt}(\epsilon \xi_i + \dot{q}_i \epsilon \tau) - \dot{q}_i \frac{d(\epsilon \tau)}{dt} \\
&= \epsilon \frac{d\xi_i}{dt} + \ddot{q}_i \epsilon \tau + \dot{q}_i \epsilon \frac{d\tau}{dt} - \dot{q}_i \epsilon \frac{d\tau}{dt} - \ddot{q}_i \epsilon \tau \\
&= \epsilon \frac{d\xi_i}{dt}
\end{align*}
\end{block}

\pause

\begin{alertblock}{关键结果}
$$\delta \dot{q}_i = \epsilon \left(\frac{\partial \xi_i}{\partial q_j}\dot{q}_j + \frac{\partial \xi_i}{\partial t}\right)$$
\end{alertblock}
\end{frame}

\begin{frame}
\frametitle{A.4 对称性条件}
\begin{block}{正确表述}
\textbf{对称性}意味着拉格朗日量至多相差全导数：

$$\delta L = L(q', \dot{q}', t') - L(q, \dot{q}, t) = \epsilon \frac{dF(q, t)}{dt}$$

其中 $F(q, t)$ 是某个函数。
\end{block}

\pause

\begin{alertblock}{重要说明}
这意味着作用量至多相差边界项：
$$\delta S = \epsilon[F(q(t_2), t_2) - F(q(t_1), t_1)]$$

由于变分时端点固定，边界项的变分为零，因此运动方程不变。
\end{alertblock}

\pause

\begin{block}{对称性条件}
$$\frac{\partial L}{\partial q_i}\xi_i + \frac{\partial L}{\partial \dot{q}_i}\frac{d\xi_i}{dt} + \left(\frac{\partial L}{\partial t} + \frac{dL}{dt}\right)\tau = \frac{dF}{dt}$$
\end{block}
\end{frame}

\begin{frame}
\frametitle{A.5 利用运动方程}
\begin{block}{代入欧拉-拉格朗日方程}
使用 $\frac{\partial L}{\partial q_i} = \frac{d}{dt}\left(\frac{\partial L}{\partial \dot{q}_i}\right)$：
\begin{align*}
&\frac{d}{dt}\left(\frac{\partial L}{\partial \dot{q}_i}\right)\xi_i + \frac{\partial L}{\partial \dot{q}_i}\frac{d\xi_i}{dt} + \left(\frac{\partial L}{\partial t} + \frac{dL}{dt}\right)\tau \\
&= \frac{d}{dt}\left(\frac{\partial L}{\partial \dot{q}_i}\xi_i\right) + \left(\frac{\partial L}{\partial t} + \frac{dL}{dt}\right)\tau
\end{align*}
\end{block}

\pause

\begin{block}{哈密顿量}
注意到：
$$\frac{dL}{dt} = \frac{\partial L}{\partial q_i}\dot{q}_i + \frac{\partial L}{\partial \dot{q}_i}\ddot{q}_i + \frac{\partial L}{\partial t}$$
使用欧拉-拉格朗日方程 $\frac{\partial L}{\partial q_i} = \frac{d}{dt}\left(\frac{\partial L}{\partial \dot{q}_i}\right) = \dot{p}_i$：
$$\frac{dL}{dt} = \dot{p}_i\dot{q}_i + p_i\ddot{q}_i + \frac{\partial L}{\partial t} = \frac{d}{dt}(p_i\dot{q}_i) + \frac{\partial L}{\partial t}$$
若 $\frac{\partial L}{\partial t} = 0$，则：
$$\frac{dH}{dt} = \frac{d}{dt}(p_i\dot{q}_i - L) = -\frac{dL}{dt}$$
\end{block}
\end{frame}

\begin{frame}
\frametitle{A.6 推导守恒量}
\begin{block}{整理对称性条件}
从对称性条件：
$$\frac{d}{dt}\left(p_i \xi_i\right) - \frac{dH}{dt}\tau = \frac{dF}{dt}$$
\end{block}

\pause

\begin{block}{定义守恒量}
重新整理得：
$$\frac{d}{dt}\left(p_i \xi_i - H\tau - F\right) = 0$$
\end{block}

\pause

\begin{alertblock}{诺特守恒量}
定义诺特守恒量：
$$Q = \sum_{i=1}^n p_i \xi_i - H\tau - F$$
则沿着运动轨迹：
$$\frac{dQ}{dt} = 0 \quad \Rightarrow \quad Q = \text{常数}$$
\end{alertblock}
\end{frame}

\begin{frame}
\frametitle{A.7 诺特定理的结论}
\begin{block}{推导过程}
由对称性条件 $\delta L = \epsilon \frac{dF}{dt}$，经过推导得到：

$$\frac{d}{dt}\left[\sum_i p_i\xi_i - H\tau + F\right] = \frac{dF}{dt}$$

因此：

$$\boxed{\frac{d}{dt}\left[\sum_i p_i\xi_i - H\tau\right] = 0}$$
\end{block}

\pause

\begin{theorem}[诺特定理]
\textbf{守恒量（诺特荷）：}
$$\boxed{Q = \sum_i p_i\xi_i - H\tau}$$

其中：
\begin{itemize}
\item $p_i = \frac{\partial L}{\partial \dot{q}_i}$ 是广义动量
\item $H = \sum_i p_i\dot{q}_i - L$ 是哈密顿量
\end{itemize}

\textbf{注：}$F$ 项不影响守恒性
\end{theorem}
\end{frame}

\begin{frame}
\frametitle{A.8 验证：空间平移对称性 → 动量守恒}
\begin{block}{空间平移变换}
考虑沿 $x$ 方向的平移：
$$t \to t, \quad \vec{r} \to \vec{r} + \epsilon\hat{x}$$
即：$\tau = 0$，$\xi_x = 1$，$\xi_y = \xi_z = 0$
\end{block}

\pause

\begin{block}{对称性条件}
若系统具有空间平移对称性（$L$ 不显含 $x$），则：
$$\frac{\partial L}{\partial x} = 0$$
对称性条件满足。
\end{block}

\pause

\begin{block}{守恒量}
应用诺特定理（$F=0$）：
$$Q = p_x \cdot 1 - H \cdot 0 = p_x$$
\end{block}

\pause

\begin{alertblock}{结论}
$$\boxed{\frac{dp_x}{dt} = 0 \quad \Rightarrow \quad \text{动量} p_x \text{守恒}}$$
\end{alertblock}
\end{frame}

\begin{frame}
\frametitle{A.9 验证：时间平移对称性 → 能量守恒}
\begin{block}{时间平移变换}
考虑纯时间平移：
$$t \to t + \epsilon, \quad q_i \to q_i$$
即：$\tau = 1$，$\xi_i = 0$
\end{block}

\pause

\begin{block}{对称性条件}
若系统具有时间平移对称性（$L$ 不显含 $t$），则：
$$\frac{\partial L}{\partial t} = 0$$
对称性条件满足。
\end{block}

\pause

\begin{block}{守恒量}
应用诺特定理（$F=0$）：
$$Q = \sum_i p_i \cdot 0 - H \cdot 1 = -H$$
\end{block}

\pause

\begin{alertblock}{结论}
$$\boxed{\frac{dH}{dt} = 0 \quad \Rightarrow \quad \text{能量} H \text{守恒}}$$
\end{alertblock}
\end{frame}

\section{附录B：场论中的诺特定理与守恒流}

\begin{frame}
\frametitle{B.1 场论的基本框架}
\begin{block}{从粒子到场}
\begin{itemize}
\item 点粒子：$q_i(t)$ - 有限自由度
\item 场论：$\phi^a(x^\mu)$ - 无限自由度
\item 时空坐标：$x^\mu = (t, \vec{x})$，$\mu = 0,1,2,3$
\end{itemize}
\end{block}

\pause

\begin{block}{拉格朗日密度}
作用量写为：
$$S = \int d^4x \, \mathcal{L}(\phi^a, \partial_\mu\phi^a)$$
其中 $\mathcal{L}$ 是拉格朗日密度。
\end{block}

\pause

\begin{block}{类比关系}
\begin{center}
\begin{tabular}{c|c}
点粒子 & 场论 \\
\hline
$q_i(t)$ & $\phi^a(x^\mu)$ \\
$L$ & $\mathcal{L}$ \\
$\int dt$ & $\int d^4x$ \\
$\dot{q}_i$ & $\partial_\mu\phi^a$
\end{tabular}
\end{center}
\end{block}
\end{frame}

\begin{frame}
\frametitle{B.2 场的运动方程}
\begin{block}{变分原理}
要求作用量 $S$ 取极值：$\delta S = 0$
\end{block}

\pause

\begin{block}{欧拉-拉格朗日方程（场论版本）}
$$\boxed{\partial_\mu\left(\frac{\partial \mathcal{L}}{\partial(\partial_\mu\phi^a)}\right) - \frac{\partial \mathcal{L}}{\partial \phi^a} = 0}$$
\end{block}

\pause

\begin{block}{例子：实标量场}
$$\mathcal{L} = \frac{1}{2}\partial_\mu\phi\partial^\mu\phi - \frac{1}{2}m^2\phi^2$$
运动方程：
$$\partial_\mu\partial^\mu\phi + m^2\phi = 0 \quad \text{（Klein-Gordon方程）}$$
\end{block}
\end{frame}

\begin{frame}
\frametitle{B.3 场的对称变换}
\begin{block}{无穷小变换}
考虑场的无穷小变换（依赖参数 $\epsilon^a$）：
$$\phi^a(x) \to \phi^a(x) + \epsilon^b \Delta_b\phi^a(x)$$
其中 $\Delta_b\phi^a$ 是变换生成元。
\end{block}

\pause

\begin{block}{坐标变换}
同时考虑时空坐标的变换：
$$x^\mu \to x^\mu + \epsilon^a \xi^\mu_a(x)$$
\end{block}

\pause

\begin{block}{场的全变分}
沿着轨迹，场的总变化：
$$\delta\phi^a = \epsilon^b\Delta_b\phi^a + (\partial_\mu\phi^a)\epsilon^b\xi^\mu_b$$
第一项：同点的场值改变 \\
第二项：坐标点改变引起的场值改变
\end{block}
\end{frame}

\begin{frame}
\frametitle{B.4 拉格朗日密度的变化}
\begin{block}{变分}
拉格朗日密度的变化：
$$\delta\mathcal{L} = \frac{\partial\mathcal{L}}{\partial\phi^a}\delta\phi^a + \frac{\partial\mathcal{L}}{\partial(\partial_\mu\phi^a)}\delta(\partial_\mu\phi^a)$$
\end{block}

\pause

\begin{block}{利用运动方程}
使用欧拉-拉格朗日方程：
$$\delta\mathcal{L} = \partial_\mu\left[\frac{\partial\mathcal{L}}{\partial(\partial_\mu\phi^a)}\delta\phi^a\right]$$
\end{block}

\pause

\begin{block}{考虑坐标变换}
完整的变化还包括时空体积元的变化：
$$\delta(d^4x \, \mathcal{L}) = d^4x \left[\delta\mathcal{L} + \partial_\mu(\mathcal{L}\epsilon^a\xi^\mu_a)\right]$$
\end{block}
\end{frame}

\begin{frame}
\frametitle{B.5 对称性条件}
\textbf{对称性}意味着拉格朗日密度至多相差四维散度：

$$\boxed{\delta\mathcal{L} = \epsilon\partial_\mu K^\mu}$$

其中 $K^\mu$ 是某个四维矢量。

\vspace{0.5em}

\textbf{说明}：积分后作用量相差边界项
$$\delta S = \epsilon\int_{\partial V} K^\mu dS_\mu$$
在场于无穷远处衰减时，此项为零。

\pause

\begin{alertblock}{关键等式}
$$\partial_\mu\left[\frac{\partial\mathcal{L}}{\partial(\partial_\mu\phi^b)}\delta\phi^b + \mathcal{L}\epsilon^a\xi^\mu_a - \epsilon^a K^\mu_a\right] = 0$$
\end{alertblock}
\end{frame}

\begin{frame}
\frametitle{B.5补充：为何全导数/散度可以忽略}

\textbf{一维情况（点粒子）：}
$$L' = L + \frac{dF}{dt} \quad \Rightarrow \quad S' = S + [F(t_2) - F(t_1)]$$

端点固定 $\Rightarrow$ 边界项对 $\delta S = 0$ 无影响

\vspace{1em}

\textbf{四维情况（场论）：}
$$\mathcal{L}' = \mathcal{L} + \partial_\mu K^\mu \quad \Rightarrow \quad S' = S + \int_{\partial V} K^\mu dS_\mu$$

场在无穷远衰减 $\Rightarrow$ 边界项为零

\vspace{1em}

\begin{alertblock}{核心思想}
运动方程只依赖于作用量的\textbf{变分}，边界项在变分时消失！
\end{alertblock}
\end{frame}

\begin{frame}
\frametitle{B.6 诺特流的定义}
\begin{block}{守恒流}
定义四维流矢量（诺特流）：
$$\boxed{j^\mu_a = \frac{\partial\mathcal{L}}{\partial(\partial_\mu\phi^b)}\Delta_a\phi^b + \mathcal{L}\xi^\mu_a - K^\mu_a}$$
\end{block}

\pause

\begin{theorem}[场论中的诺特定理]
若系统在变换 $\phi^a \to \phi^a + \epsilon^b\Delta_b\phi^a$ 下具有对称性，则存在守恒流 $j^\mu_a$ 满足：
$$\boxed{\partial_\mu j^\mu_a = 0}$$
这是连续性方程。
\end{theorem}

\pause

\begin{block}{物理意义}
\begin{itemize}
\item $j^0_a$：守恒荷密度
\item $\vec{j}_a$：守恒荷流密度
\item $\partial_\mu j^\mu = 0$：局域守恒律
\end{itemize}
\end{block}
\end{frame}

\begin{frame}
\frametitle{B.7 守恒荷}
\begin{block}{积分守恒量}
从连续性方程 $\partial_\mu j^\mu_a = 0$，定义守恒荷：
$$Q_a = \int d^3x \, j^0_a$$
\end{block}

\pause

\begin{block}{时间守恒}
对 $Q_a$ 求时间导数：
\begin{align*}
\frac{dQ_a}{dt} &= \int d^3x \, \frac{\partial j^0_a}{\partial t} = -\int d^3x \, \nabla \cdot \vec{j}_a \\
&= -\int_{\text{边界}} \vec{j}_a \cdot d\vec{S} = 0
\end{align*}
假设场在无穷远处趋于零。
\end{block}

\pause

\begin{alertblock}{守恒律}
$$\boxed{\frac{dQ_a}{dt} = 0 \quad \Rightarrow \quad Q_a = \text{常数}}$$
这是全局守恒量。
\end{alertblock}
\end{frame}

\begin{frame}
\frametitle{B.8 例子：复标量场的 U(1) 对称性（I）}
\begin{block}{拉格朗日密度}
$$\mathcal{L} = \partial_\mu\phi^*\partial^\mu\phi - m^2\phi^*\phi$$
其中 $\phi$ 是复标量场。
\end{block}

\pause

\begin{block}{U(1) 变换}
相位变换：
$$\phi \to e^{i\alpha}\phi, \quad \phi^* \to e^{-i\alpha}\phi^*$$
无穷小形式（$\alpha = \epsilon$）：
$$\delta\phi = i\epsilon\phi, \quad \delta\phi^* = -i\epsilon\phi^*$$
\end{block}

\pause

\begin{block}{对称性验证}
$$\delta\mathcal{L} = \partial_\mu(i\epsilon\phi^*)\partial^\mu\phi + \partial_\mu\phi^*\partial^\mu(i\epsilon\phi) - m^2(i\epsilon\phi^*\phi - i\epsilon\phi^*\phi) = 0$$
拉格朗日密度严格不变！
\end{block}
\end{frame}

\begin{frame}
\frametitle{B.9 例子：复标量场的 U(1) 对称性（II）}
\begin{block}{诺特流}
纯内对称（无时空变换），$\xi^\mu = 0$，$K^\mu = 0$：
$$j^\mu = \frac{\partial\mathcal{L}}{\partial(\partial_\mu\phi)}\delta\phi + \frac{\partial\mathcal{L}}{\partial(\partial_\mu\phi^*)}\delta\phi^*$$
\end{block}

\pause

\begin{block}{计算}
$$\frac{\partial\mathcal{L}}{\partial(\partial_\mu\phi)} = \partial^\mu\phi^*, \quad \frac{\partial\mathcal{L}}{\partial(\partial_\mu\phi^*)} = \partial^\mu\phi$$
$$j^\mu = \partial^\mu\phi^* \cdot (i\epsilon\phi) + \partial^\mu\phi \cdot (-i\epsilon\phi^*) = i\epsilon(\phi^*\partial^\mu\phi - \phi\partial^\mu\phi^*)$$
\end{block}

\pause

\begin{alertblock}{守恒流}
$$\boxed{j^\mu = i(\phi^*\partial^\mu\phi - \phi\partial^\mu\phi^*)}$$
这正是电磁学中的电流密度形式！守恒荷 $Q = \int d^3x \, j^0$ 对应电荷。
\end{alertblock}
\end{frame}

\begin{frame}
\frametitle{B.10 能动张量与四动量（I）}
\begin{block}{时空平移对称性}
考虑时空平移：
$$x^\mu \to x^\mu + \epsilon^\mu$$
即：$\xi^\mu_\nu = \delta^\mu_\nu$，场值不变 $\Delta_\nu\phi^a = 0$
\end{block}

\pause

\begin{block}{对称性条件}
若 $\mathcal{L}$ 不显含时空坐标 $x^\mu$，则具有平移对称性。
\end{block}

\pause

\begin{block}{诺特流 - 能动张量}
$$j^\mu_\nu = \frac{\partial\mathcal{L}}{\partial(\partial_\mu\phi^a)} \cdot 0 + \mathcal{L}\delta^\mu_\nu = \mathcal{L}\delta^\mu_\nu$$
这不够完整！需要考虑场的变化：
$$T^\mu_\nu = \frac{\partial\mathcal{L}}{\partial(\partial_\mu\phi^a)}\partial_\nu\phi^a - \mathcal{L}\delta^\mu_\nu$$
\end{block}
\end{frame}

\begin{frame}
\frametitle{B.11 能动张量与四动量（II）}
\begin{block}{能动张量}
$$\boxed{T^\mu_\nu = \frac{\partial\mathcal{L}}{\partial(\partial_\mu\phi^a)}\partial_\nu\phi^a - \mathcal{L}\delta^\mu_\nu}$$
守恒条件：
$$\partial_\mu T^\mu_\nu = 0$$
\end{block}

\pause

\begin{block}{四动量}
守恒荷（四动量）：
$$P^\nu = \int d^3x \, T^0_\nu$$
\begin{itemize}
\item $P^0 = \int d^3x \, T^{00}$：能量
\item $\vec{P} = \int d^3x \, \vec{T}^0$：动量
\end{itemize}
\end{block}

\pause

\begin{alertblock}{守恒律}
$$\frac{dP^\nu}{dt} = 0 \quad \Rightarrow \quad \text{四动量守恒}$$
\end{alertblock}
\end{frame}

\begin{frame}
\frametitle{B.12 对称性与守恒流对应表}
\begin{center}
\begin{tabular}{|c|c|c|}
\hline
\textbf{对称性} & \textbf{守恒流} & \textbf{守恒荷} \\
\hline
时空平移 & 能动张量 $T^\mu_\nu$ & 四动量 $P^\nu$ \\
\hline
空间旋转 & 角动量流 $M^\mu_{\rho\sigma}$ & 角动量 $J^{\rho\sigma}$ \\
\hline
U(1) 相位 & 电流 $j^\mu$ & 电荷 $Q$ \\
\hline
SU(2) 同位旋 & 同位旋流 $I^\mu_a$ & 同位旋 $I_a$ \\
\hline
SU(3) 色 & 色流 $j^\mu_c$ & 色荷 $Q_c$ \\
\hline
\end{tabular}
\end{center}

\vspace{1em}

\begin{block}{统一框架}
所有守恒律都源于对称性，通过诺特定理建立联系：
$$\text{对称性} \xrightarrow{\text{诺特定理}} \text{守恒流} \xrightarrow{\text{积分}} \text{守恒荷}$$
\end{block}
\end{frame}

\begin{frame}
\frametitle{B.13 场论 vs 点粒子力学对比}
\begin{center}
\begin{tabular}{|c|c|c|}
\hline
\textbf{概念} & \textbf{点粒子} & \textbf{场论} \\
\hline
自由度 & $q_i(t)$ & $\phi^a(x^\mu)$ \\
\hline
拉格朗日 & $L$ & $\mathcal{L}$ \\
\hline
作用量 & $\int dt \, L$ & $\int d^4x \, \mathcal{L}$ \\
\hline
运动方程 & $\frac{d}{dt}\frac{\partial L}{\partial\dot{q}_i} - \frac{\partial L}{\partial q_i} = 0$ & $\partial_\mu\frac{\partial\mathcal{L}}{\partial(\partial_\mu\phi^a)} - \frac{\partial\mathcal{L}}{\partial\phi^a} = 0$ \\
\hline
守恒量 & $Q = p_i\xi_i - H\tau$ & $Q = \int d^3x \, j^0$ \\
\hline
守恒律 & $\frac{dQ}{dt} = 0$ & $\partial_\mu j^\mu = 0$ \\
\hline
\end{tabular}
\end{center}

\vspace{1em}

\begin{alertblock}{核心思想}
场论是点粒子力学在\textcolor{red}{连续极限}下的推广，诺特定理的形式更加\textcolor{red}{局域化}。
\end{alertblock}
\end{frame}

\begin{frame}
\frametitle{B.14 从全局到局域对称性}
\begin{block}{全局对称性}
变换参数 $\epsilon$ 与时空坐标无关：
$$\phi(x) \to \phi(x) + \epsilon \Delta\phi(x)$$
导致守恒流 $j^\mu$，$\partial_\mu j^\mu = 0$
\end{block}

\pause

\begin{block}{局域对称性（规范对称性）}
变换参数依赖时空：$\epsilon \to \epsilon(x)$
$$\phi(x) \to \phi(x) + \epsilon(x)\Delta\phi(x)$$
为保持对称性，必须引入\textcolor{red}{规范场} $A^\mu(x)$！
\end{block}

\pause

\begin{block}{协变导数}
$$\partial_\mu \to D_\mu = \partial_\mu - iqA_\mu$$
\end{block}

\pause

\begin{alertblock}{深刻意义}
局域对称性要求 $\Rightarrow$ 必须引入相互作用！\\
这是现代规范场论的基础（电磁、弱、强相互作用）。
\end{alertblock}
\end{frame}

\begin{frame}
\frametitle{B.15 小结}
\begin{block}{场论中的诺特定理}
\begin{enumerate}
\item 对称变换 $\to$ 守恒流 $j^\mu$
\item 连续性方程：$\partial_\mu j^\mu = 0$
\item 守恒荷：$Q = \int d^3x \, j^0$，$\frac{dQ}{dt} = 0$
\end{enumerate}
\end{block}

\pause

\begin{block}{关键要点}
\begin{itemize}
\item 场论守恒律是\textcolor{red}{局域的}（每个时空点都满足）
\item 能动张量 $T^\mu_\nu$ 描述能量-动量分布
\item 内对称性（如 U(1)）对应守恒荷（如电荷）
\item 时空对称性对应能量-动量守恒
\end{itemize}
\end{block}

\pause

\begin{alertblock}{现代物理}
场论中的诺特定理是量子场论、粒子物理标准模型的\textcolor{red}{数学基础}。
\end{alertblock}
\end{frame}

\end{document}